\chapter{Formularze}\label{dod:formularze}

\begin{blok}
Wydrukuj ten dodatek osobno dla każdej misji (jednostronnie). Wypełniaj formularze długopisem. Każdy formularz ma:
\begin{itemize}
  \item nagłówek z danymi misji i osobą wypełniającą,
  \item kroki w formie poleceń: \emph{wejdź na \dots{} i sprawdź \dots{}},
  \item pole „Wartość / uwagi” na odczyty, numery i nazwiska,
  \item miejsce na komentarze i podpisy.
\end{itemize}
Gdy krok jest niewykonalny, nie zaznaczaj go -- opisz odstępstwo w komentarzach. Kierownik misji decyduje, czy kontynuować.

Po misji zeskanuj wszystkie formularze i przechowuj skany w wewnętrznej dokumentacji koła. Wypełnione formularze zawierają dane osobowe (nazwiska, telefony) -- nie publikuj ich. Do publicznego raportu przenieś tylko dane techniczne i wnioski.

\medskip
\renewcommand{\arraystretch}{1.25}
\begin{tabularx}{\linewidth}{@{}P{12mm}LP{38mm}@{}}
\toprule
\textbf{Kod} & \textbf{Formularz} & \textbf{Kiedy}\\\midrule
F0 & Karta misji & cały projekt\\
F1 & Role w zespole i macierz odpowiedzialności & T$-$16 tyg.\\
F2 & Planowanie misji & T$-$16 do T$-$4 tyg.\\
F3 & Formalności & T$-$4 do T$-$2 tyg.\\
F4 & Dziennik predykcji & od T$-$7 dni\\
F5 & Raport z testów & przed FRR\\
F6 & Przegląd gotowości do startu (FRR) & T$-$7 do T$-$3 dni\\
F7 & Konfiguracja trackera i serwisów & T$-$2 dni\\
F8 & Weryfikacja i logistyka & T$-$1 dzień\\
F9 & Stanowisko startowe & T0\\
F10 & Tankowanie & T0\\
F11 & Start & T0\\
F12 & Lot i pościg (dziennik) & T0\\
F13 & Odzyskanie & T0\\
F14 & Po misji & do T+7 dni\\
\bottomrule
\end{tabularx}
\end{blok}

% ===========================================================================
\formularz[kierownik misji]{F0}{Karta misji (1/2)}{cały projekt}
\renewcommand{\arraystretch}{1.0}
\small
\begin{tabularx}{\linewidth}{@{}|L|L|@{}}\hline
\multicolumn{2}{|l|}{\cellcolor{sdlight}\rule{0pt}{4mm}\textbf{Misja}}\\\hline
\etykieta{Nazwa misji} & \etykieta{Kierownik misji (telefon)}\\\hline
\multicolumn{2}{|p{\dimexpr\linewidth-2\tabcolsep\relax}|}{\etykieta[40mm]{Cel misji (mission statement): co chcemy osiągnąć, po czym poznamy sukces}}\\\hline
\multicolumn{2}{|p{\dimexpr\linewidth-2\tabcolsep\relax}|}{\etykieta[18mm]{Ładunki badawcze (nazwa, zespół, opiekun)}}\\\hline
\multicolumn{2}{|l|}{\cellcolor{sdlight}\rule{0pt}{4mm}\textbf{Terminy i miejsca (zgodne ze zgłoszeniem PAŻP)}}\\\hline
\etykieta[11mm]{Scenariusz A: miejsce, współrzędne} & \etykieta[11mm]{Scenariusz A: data, godz. (UTC), termin zapasowy}\\\hline
\etykieta[11mm]{Scenariusz B: miejsce, współrzędne} & \etykieta[11mm]{Scenariusz B: data, godz. (UTC), termin zapasowy}\\\hline
\etykieta{Numery zgłoszeń PAŻP} & \etykieta{NOTAM (numer, sprawdzony w AIS: data)}\\\hline
\etykieta{Telefon do służby ruchu lotniczego} & \etykieta{Warunki zgody (np. telefon 15 min przed startem)}\\\hline
\etykieta{Ubezpieczenie (nr polisy)} & \etykieta{Zgoda właściciela terenu A / B}\\\hline
\end{tabularx}
\normalsize
\podpisy

\formularz[kierownik misji]{F0}{Karta misji (2/2)}{cały projekt}
\renewcommand{\arraystretch}{1.0}
\small
\begin{tabularx}{\linewidth}{@{}|L|L|@{}}\hline
\multicolumn{2}{|l|}{\cellcolor{sdlight}\rule{0pt}{4mm}\textbf{Łączność i lokalizacja}}\\\hline
\etykieta{Znak wywoławczy i SSID APRS} & \etykieta{Właściciel znaku, data zgody na użycie}\\\hline
\etykieta{Payload ID (RTTY), częstotliwości} & \etykieta{Kanały PMR: główny / zapasowy}\\\hline
\etykieta{Identyfikator SPOT} & \etykieta{Nazwy znaczników Bluetooth (Apple / Google)}\\\hline
\etykieta{Wersja kodu HABduino (release / commit)} & \etykieta{Wersja kodu OBC (release / commit)}\\\hline
\etykieta{Wersja kodu ładunku (release / commit)} & \etykieta{Znane błędy i procedury obejścia}\\\hline
\multicolumn{2}{|l|}{\cellcolor{sdlight}\rule{0pt}{4mm}\textbf{Zestaw lotny}}\\\hline
\etykieta{Masa ładunku bez powłoki [g] / klasa balonu} & \etykieta{Masa pod spadochronem [g]}\\\hline
\etykieta{Powłoka (model, masa)} & \etykieta{Spadochron (kolor, średnica, $C_d$, prędkość opadania)}\\\hline
\etykieta{Prędkość wznoszenia [m/s]} & \etykieta{Siła nośna szyjki = masa obciążnika [g]}\\\hline
\etykieta{Objętość gazu [m³]} & \etykieta{Butle: ilość, pojemność, czystość, miejsce odbioru}\\\hline
\etykieta{Przewidywana wysokość pęknięcia [m]} & \etykieta{Przewidywany czas lotu [min]}\\\hline
\etykieta{Długość zapasu linki odcinacza [m]} & \etykieta{Progi odcinacza (czas, wysokość, obszar)}\\\hline
\end{tabularx}
\normalsize
\podpisy

% ===========================================================================
\formularz[kierownik misji]{F1}{Role w zespole}{T$-$16 tyg.}
Jedna osoba może mieć kilka ról, z wyjątkiem: kierowca nie obsługuje śledzenia; fotograf nie ma innych zadań przy starcie; kierownik startu nie tankuje.

\renewcommand{\arraystretch}{1.55}
\small
\begin{tabularx}{\linewidth}{@{}|P{11mm}|P{48mm}|L|P{30mm}|@{}}\hline
\textbf{Kod} & \textbf{Rola i główne zadania} & \textbf{Imię i nazwisko} & \textbf{Telefon}\\\hline
KM & Kierownik misji: decyzje „leci / nie leci”, odpowiedzialność za całość & & \\\hline
KS & Kierownik startu: komendy na stanowisku, prowadzi formularze F9--F11 & & \\\hline
FL & Formalności i łączność z PAŻP, właściciel terenu & & \\\hline
OT & Operator tankowania: butla, reduktor, szyjka & & \\\hline
PM & Pomocnicy (2--4 os.): folia, powłoka, start & & \\\hline
OP & Operator platformy: gondola, tracker, OBC, odcinacz & & \\\hline
OŁ & Opiekun ładunku badawczego & & \\\hline
OS & Operator śledzenia: predykcja, mapy, nawigacja & & \\\hline
KR & Kierowca (każdego samochodu) & & \\\hline
FD & Fotograf / dokumentalista & & \\\hline
LG & Logistyka i zakupy: gaz, materiały, transport & & \\\hline
BZ & Bezpieczeństwo: butla, pierwsza pomoc, widzowie & & \\\hline
\end{tabularx}
\normalsize
\komentarze[2]
\podpisy

% ---------------------------------------------------------------- RACI
\clearpage
\begin{landscape}
\markright{F1 Macierz odpowiedzialności}
\noindent{\sffamily\Large\bfseries\color{sdnavy} F1\quad Macierz odpowiedzialności}\par\smallskip
\noindent\textbf{R} -- wykonuje, \textbf{A} -- odpowiada i zatwierdza (jedna osoba), \textbf{C} -- konsultowany, \textbf{I} -- informowany. Puste pole: rola nie bierze udziału.\par\smallskip
\renewcommand{\arraystretch}{1.12}
\footnotesize
\begin{tabularx}{\linewidth}{@{}|L|*{12}{C{9mm}|}@{}}\hline
\rowcolor{sdlight}\textbf{Czynność} & \textbf{KM} & \textbf{KS} & \textbf{FL} & \textbf{OT} & \textbf{PM} & \textbf{OP} & \textbf{OŁ} & \textbf{OS} & \textbf{KR} & \textbf{FD} & \textbf{LG} & \textbf{BZ}\\\hline
\multicolumn{13}{|l|}{\cellcolor{sdlight!50}\itshape Przygotowania}\\\hline
Cel misji, wymagania ładunku & A & & & & & C & R & & & & & \\\hline
Budżet masy, obliczenie gazu & A & & & C & & R & C & & & & & \\\hline
Zgłoszenia PAŻP (dwa scenariusze) & A & & R & & & & & C & & & & \\\hline
Ubezpieczenie, zgoda właściciela terenu & A & & R & & & & & & & & & I\\\hline
Zamówienie gazu i materiałów & A & & & C & & & & & & & R & \\\hline
Testy platformy / testy ładunku & A & & & & & R & R & & & C & & \\\hline
Predykcje, dziennik predykcji & A & & I & & & & & R & I & & & \\\hline
Konfiguracja trackera i serwisów (T$-$2) & A & & & & & R & & C & & & & \\\hline
Przegląd gotowości, decyzja & A/R & C & C & C & & C & C & C & I & I & C & C\\\hline
Pakowanie i transport sprzętu & A & I & & C & & C & C & & C & & R & \\\hline
\multicolumn{13}{|l|}{\cellcolor{sdlight!50}\itshape Dzień startu}\\\hline
Rozłożenie i zabezpieczenie stanowiska & I & A & & C & R & & & & & I & & C\\\hline
Montaż reduktora, tankowanie & I & A & & R & C & & & & & I & & C\\\hline
Zamknięcie szyjki, mocowanie olinowania & I & A & & R & C & C & & & & I & & \\\hline
Gondola platformy, uzbrojenie odcinacza & I & A & & & & R & C & C & & I & & \\\hline
Włączenie i zamknięcie ładunku & I & A & & & & C & R & & & I & & \\\hline
Kontrola telemetrii przed startem & I & A & & & & C & & R & & & & \\\hline
Telefon do służby ruchu lotniczego & A & I & R & & & & & & & & & \\\hline
Decyzja o starcie & A/R & C & C & & & C & C & C & & & & C\\\hline
Start (puszczenie zestawu) & I & A/R & & & R & C & & I & & C & & C\\\hline
Bezpieczeństwo osób postronnych & I & A & & & & & & & & & & R\\\hline
Dokumentacja zdjęciowa & A & & & & & & & & & R & & \\\hline
Śledzenie i nawigacja & A & & & & & & & R & C & & & \\\hline
Jazda samochodem pościgowym & A & & & & & & & C & R & & & \\\hline
Odzyskanie, zabezpieczenie danych & A & & & & & R & R & C & C & C & & \\\hline
Zgłoszenie końca lotu & A & & R & & & & & I & & & & \\\hline
\multicolumn{13}{|l|}{\cellcolor{sdlight!50}\itshape Po misji}\\\hline
Raport z misji, wnioski, aktualizacja podręcznika & A & C & C & C & & R & R & C & & C & C & \\\hline
\end{tabularx}
\normalsize
\end{landscape}

% ===========================================================================
\formularz[kierownik misji]{F2}{Planowanie misji}{T$-$16 do T$-$4 tyg.}
\begin{procedura}
\etap{Cel i ładunek}
\krok{Zapisz cel misji i wymagania ładunku w karcie misji (F0).}{}
\krok{Wypełnij role w zespole (F1). Przekaż macierz odpowiedzialności całemu zespołowi.}{}
\krok{Uzgodnij z zespołem badawczym interfejs: pudełko, masa, zasilanie, częstotliwości.}{\pole[2cm]{masa ład. [g]:}}
\krok{Ustal terminy przeglądów (CDR, FRR) i wpisz je do kalendarza zespołu.}{\pole[2cm]{FRR:}}
\etap{Sprzęt}
\krok{Wyjmij sprzęt z szafki. Sprawdź go według wykazu materiałów (rozdział~\ref{ch:materialy}).}{\pole[2cm]{braki:}}
\krok{Sprawdź datę ważności i stan powłoki (bez pęknięć, bez lepkości).}{}
\krok{Wypełnij budżet masy (rozdział~\ref{ch:planowanie}). Sprawdź klasę: lekki.}{\pole[2cm]{razem [g]:}}
\krok{Wejdź na \www{sondehub.org/calc} i oblicz ilość gazu dla \SI{5}{m/s}. Wpisz wyniki do F0.}{\pole[2cm]{gaz [m³]:}}
\etap{Wstępny plan lotu}
\krok{Wybierz miejsce zasadnicze (A) i zapasowe (B). Uzyskaj zgodę właścicieli terenu.}{}
\krok{Wejdź na \www{predict.sondehub.org} i sprawdź typowe trasy dla A i B w tej porze roku.}{}
\krok{Wybierz znak wywoławczy: klubowy SimLE, członka zespołu albo udostępniony. Uzyskaj zgodę właściciela (rozdział~\ref{ch:prawo}).}{\pole[2cm]{znak:}}
\end{procedura}
\komentarze[3]
\podpisy

% ===========================================================================
\formularz[formalności (FL)]{F3}{Formalności}{T$-$4 do T$-$2 tyg.}
\begin{procedura}
\etap{Zgłoszenie do PAŻP}
\krok{Wejdź na \www{www.pansa.pl/asm1/formularze-zgloszeniowe} i pobierz aktualny formularz sondowania atmosfery.}{}
\krok{Wypełnij formularz dla scenariusza A: termin zasadniczy i zapasowy.}{}
\krok{Wypełnij formularz dla scenariusza B: termin zasadniczy i zapasowy.}{}
\krok{Wyślij oba formularze na \href{mailto:asm1@pansa.pl}{\nolinkurl{asm1@pansa.pl}}, co najmniej 14 dni przed lotem.}{\pole[2cm]{data:}}
\krok{Zapisz potwierdzenie i numer telefonu do służby ruchu lotniczego w F0.}{\pole[2cm]{tel.:}}
\etap{Ubezpieczenie i zgody}
\krok{Wykup / potwierdź ubezpieczenie OC na lot.}{\pole[2cm]{nr polisy:}}
\krok{Potwierdź na piśmie zgodę właściciela terenu A i B.}{}
\etap{Zakupy}
\krok{Zamów gaz z zapasem co najmniej 30\,\%. Potwierdź czystość helu i termin odbioru.}{\pole[2cm]{butle:}}
\krok{Kup materiały jednorazowe: powłoka, linka, folia, baterie litowe, taśma.}{}
\krok{Aktywuj subskrypcję SPOT na miesiąc lotu.}{}
\end{procedura}
\komentarze[3]
\podpisy

% ===========================================================================
\clearpage
\begin{landscape}
\formularz[operator śledzenia]{F4}{Dziennik predykcji}{od T$-$7 dni}
\vspace{-2mm}
\noindent Codziennie wejdź na \www{predict.sondehub.org} i wpisz parametry z karty misji. Zapisz wynik dla każdego scenariusza. Decyzja: \textbf{T} -- start możliwy, \textbf{N} -- niemożliwy, \textbf{?} -- obserwuj.\par\smallskip
\renewcommand{\arraystretch}{1.4}
\footnotesize
\begin{tabularx}{\linewidth}{@{}|P{16mm}|P{8mm}|P{14mm}|P{12mm}|P{12mm}|P{30mm}|P{12mm}|P{20mm}|P{22mm}|L|P{8mm}|@{}}\hline
\rowcolor{sdlight}\textbf{Data} & \textbf{Scen.} & \textbf{Godz. startu} & \textbf{Wzn. [m/s]} & \textbf{Pękn. [km]} & \textbf{Lądowanie (współrzędne, miejscowość)} & \textbf{Odl. [km]} & \textbf{Wiatr przy ziemi} & \textbf{Strefy / NOTAM na trasie} & \textbf{Uwagi (teren, pogoda, chmury)} & \textbf{T/N}\\\hline
 & A & & & & & & & & & \\\hline
 & B & & & & & & & & & \\\hline
 & A & & & & & & & & & \\\hline
 & B & & & & & & & & & \\\hline
 & A & & & & & & & & & \\\hline
 & B & & & & & & & & & \\\hline
 & A & & & & & & & & & \\\hline
 & B & & & & & & & & & \\\hline
 & A & & & & & & & & & \\\hline
 & B & & & & & & & & & \\\hline
\end{tabularx}
\normalsize
\podpisy
\end{landscape}

% ===========================================================================
\formularz[operator platformy, opiekun ładunku]{F5}{Raport z testów}{przed FRR}
Wykonaj każdy test dla platformy i dla każdego ładunku (rozdział~\ref{ch:planowanie}). Zapisz wynik. Test niezaliczony powtórz po naprawie w nowym wierszu.

\renewcommand{\arraystretch}{2.0}
\small
\begin{tabularx}{\linewidth}{@{}|P{30mm}|P{18mm}|P{14mm}|L|P{12mm}|P{16mm}|@{}}\hline
\rowcolor{sdlight}\textbf{Test} & \textbf{Obiekt} & \textbf{Data} & \textbf{Wynik / pomiar} & \textbf{OK?} & \textbf{Podpis}\\\hline
Funkcjonalny & & & & & \\\hline
W zimnie (metoda, temp., czas) & & & & & \\\hline
Zasięg radia ($\geq$ 100 m) & & & & & \\\hline
Odcinacz główny (obciążenie, skręt, autotest) & & & & & \\\hline
Odcinacz zapasowy (tranzystor, prąd) & & & & & \\\hline
Czas pracy na bateriach & & & & & \\\hline
Upadek z 1 m & & & & & \\\hline
Równowaga gondoli & & & & & \\\hline
Próba generalna & & & & & \\\hline
 & & & & & \\\hline
\end{tabularx}
\normalsize
\komentarze[2]
\podpisy

% ===========================================================================
\formularz[kierownik misji]{F6}{Przegląd gotowości do startu}{T$-$7 do T$-$3 dni}
Każdy punkt wymaga dowodu (wynik testu, e-mail, zdjęcie). Brak dowodu oznacza „nie leci”.
\begin{procedura}
\etap{Formalności}
\krok{Zgłoszenia PAŻP potwierdzone dla scenariuszy A i B.}{}
\krok{Ubezpieczenie aktywne w dniu lotu.}{}
\krok{Zgoda właściciela terenu potwierdzona.}{}
\etap{Sprzęt}
\krok{Wszystkie testy z F5 zaliczone (platforma i ładunki).}{}
\krok{Masa ładunku poniżej 4 kg, każdy pakiet poniżej 2 kg.}{\pole[2cm]{[g]:}}
\krok{Kod lotny zamrożony. Wersje wpisane do F0.}{}
\krok{Gaz zamówiony w ilości z zapasem. Termin odbioru potwierdzony.}{\pole[2cm]{[m³]:}}
\krok{Wszystkie materiały z wykazu dostępne i spakowane.}{}
\etap{Zespół i plan}
\krok{Wszystkie role z F1 obsadzone. Każdy zna swoje zadania.}{}
\krok{Plan działania dla scenariuszy A, B i „brak startu” znany zespołowi.}{}
\krok{Predykcje (F4) nie wykluczają startu.}{}
\krok{Samochody, kierowcy i trasa dojazdu ustalone.}{}
\end{procedura}
\komentarze[3]
\podpisy

% ===========================================================================
\formularz[operator platformy]{F7}{Konfiguracja trackera i serwisów}{najpóźniej T$-$2 dni}
\begin{procedura}
\etap{Tracker}
\krok{Wgraj program lotny na HABduino. Zapisz wersję w F0.}{\pole[2cm]{wersja:}}
\krok{Sprawdź w programie: znak wywoławczy, SSID, payload ID, częstotliwości, tryb GPS \emph{airborne}.}{}
\krok{Wgraj program lotny na OBC. Sprawdź progi odcinacza zapasowego.}{\pole[2cm]{wersja:}}
\etap{Serwisy śledzenia}
\krok{Wejdź na \www{amateur.sondehub.org} i skonfiguruj dokument ładunku RTTY.}{}
\krok{Włącz tracker na zewnątrz, z anteną. Odejdź co najmniej 100 m.}{}
\krok{Wejdź na \www{aprs.fi} i sprawdź pozycję znaku z SSID.}{\pole[2cm]{godz.:}}
\krok{Wejdź na \www{amateur.sondehub.org} i sprawdź pakiety RTTY.}{\pole[2cm]{godz.:}}
\krok{Wejdź na stronę SPOT i sprawdź pozycję urządzenia.}{}
\krok{Otwórz aplikacje Apple i Google. Sprawdź pozycję obu znaczników.}{}
\krok{Uruchom stacje naziemne (horus-gui / dl-fldigi, Direwolf). Sprawdź, czy ich pakiety trafiają do SondeHub i \www{aprs.fi}.}{}
\krok{Wyślij zespołowi linki do śledzenia.}{}
\end{procedura}
\komentarze[3]
\podpisy

% ===========================================================================
\formularz[kierownik misji]{F8}{Weryfikacja i logistyka}{T$-$1 dzień}
\begin{procedura}
\etap{Pogoda i decyzja}
\krok{Wejdź na \www{predict.sondehub.org} i wykonaj predykcję dla A i B. Wpisz do F4.}{}
\krok{Sprawdź prognozę wiatru przy ziemi, opadów i burz.}{\pole[2cm]{wiatr:}}
\krok{Wejdź na \www{amc.pansa.pl} (AUP na jutro) i \www{ais.pansa.pl} (NOTAM). Sprawdź strefy na trasie i NOTAM misji.}{\pole[2cm]{NOTAM:}}
\krok{Kierownik misji wybiera scenariusz albo odwołuje lot.}{\pole[2cm]{scen.:}}
\etap{Sprzęt}
\krok{Sprawdź, czy tracker nadal jest widoczny na \www{aprs.fi} i SondeHub.}{}
\krok{Naładuj telefony, laptopy, powerbanki, radiotelefony (APRS i PMR).}{}
\krok{Włóż nowe baterie litowe (albo przygotuj je do włożenia na stanowisku).}{}
\krok{Spakuj sprzęt według wykazu materiałów. Zaznacz kolumnę „OK”.}{}
\krok{Odbierz gaz albo potwierdź odbiór na miejscu startu.}{\pole[2cm]{bar:}}
\etap{Logistyka}
\krok{Potwierdź z zespołem godzinę i miejsce zbiórki.}{\pole[2cm]{godz.:}}
\krok{Wydrukuj formularze F9--F14 i weź podkładki i długopisy.}{}
\end{procedura}
\komentarze[3]
\podpisy

% ===========================================================================
\formularz[kierownik startu]{F9}{Stanowisko startowe}{T0}
\begin{procedura}
\etap{Dokumentacja}
\krok{Fotograf: zaczyna dokumentację -- zespół, sprzęt, otwarte gondole, stanowisko. Zdjęcia przez cały dzień.}{}
\etap{Teren}
\krok{Sprawdź teren: brak drzew, linii energetycznych, przeszkód w kierunku wiatru.}{}
\krok{Sprawdź aktualną predykcję dla wybranego scenariusza.}{\pole[2cm]{lądow.:}}
\krok{Wejdź na \www{amc.pansa.pl} (AUP/UUP) i \www{ais.pansa.pl} (NOTAM). Sprawdź strefy aktywne dziś.}{}
\krok{Ustaw stanowisko w osi wiatru. Rozdaj radiotelefony PMR.}{\pole[2cm]{kanał:}}
\krok{Wyznacz strefę dla widzów (BZ).}{}
\etap{Podłoże}
\krok{Rozłóż plandekę. Dociśnij rogi skrzynkami.}{}
\krok{Rozłóż koc na plandece.}{}
\krok{Rozłóż folię dolną. Dociśnij rogi.}{}
\krok{Przygotuj folię górną obok stanowiska.}{}
\etap{Sprzęt}
\krok{Połóż butlę stabilnie. Zabezpiecz przed przetoczeniem.}{}
\krok{Rozłóż olinowanie w linii, z wiatrem.}{}
\krok{Włącz elektronikę platformy. Sprawdź GPS i telemetrię na mapach.}{\pole[2cm]{godz.:}}
\krok{Włącz ładunek (OŁ).}{}
\end{procedura}
\komentarze[3]
\podpisy

% ===========================================================================
\formularz[operator tankowania]{F10}{Tankowanie}{T0}
\begin{procedura}
\etap{Przygotowanie}
\krok{Wszyscy przy powłoce zakładają rękawiczki.}{}
\krok{Przykręć reduktor do butli. Pokrętło reduktora zamknięte.}{}
\krok{Otwórz zawór butli. Odczytaj $p_1$ z manometru ciśnienia w butli.}{\pole[2cm]{$p_1$ [bar]:}}
\krok{Zważ obciążnik: siła nośna szyjki z F0 (nie masa gondoli).}{\pole[2cm]{[g]:}}
\krok{Wyjmij powłokę na folię dolną. Włóż rurę w szyjkę. Uszczelnij taśmą delikatnie.}{}
\krok{Przywiąż obciążnik. Przykryj powłokę folią górną (PM trzymają brzegi).}{}
\etap{Tankowanie}
\krok{Otwórz reduktor powoli. Ustaw stały przepływ. Zapisz godzinę i wskazanie przepływu.}{\pole[2cm]{l/min:}}
\krok{Zapisuj odczyty manometru wysokiego ciśnienia co kilka minut (pole niżej).}{}
\krok{Balon unosi obciążnik: zamknij zawór. Odczekaj kilka minut. Odczytaj $p_2$.}{\pole[2cm]{$p_2$ [bar]:}}
\krok{Oblicz objętość gazu z tabeli~\ref{tab:gaz} albo wzoru.}{\pole[2cm]{V [m³]:}}
\etap{Zamknięcie}
\krok{Zdejmij taśmę wstępną. Wysuń rurę. Skręć szyjkę.}{}
\krok{Zwiąż szyjkę sznurkiem na kilka supłów. Przełóż przez szyjkę pętlę z krętlikiem.}{}
\krok{Zegnij szyjkę w górę (pętla w zgięciu). Zaciśnij kilkoma opaskami. Owiń kilkoma warstwami taśmy.}{}
\krok{Sprawdź: brak taśmy na powłoce, pętla trzyma, wiązanie szczelne.}{}
\end{procedura}
\renewcommand{\arraystretch}{1.4}
\noindent\begin{tabularx}{\linewidth}{@{}|l|*{8}{X|}@{}}\hline
\footnotesize godz. & & & & & & & & \\\hline
\footnotesize bar & & & & & & & & \\\hline
\end{tabularx}
\par\medskip\noindent{\sffamily\bfseries Przebieg tankowania, uwagi do procedury (co zmienić w rozdziale~\ref{ch:naziemny}):}
\foreach \i in {1,...,2}{\linia}\par
\podpisy

% ===========================================================================
\formularz[kierownik startu]{F11}{Start}{T0}
\begin{procedura}
\etap{Gotowość zestawu}
\krok{Pętla z krętlikiem zamocowana na szyjce. Węzły zabezpieczone taśmą.}{}
\krok{Linka odcinacza, zapas i spadochron rozłożone w linii, bez splątań.}{}
\krok{Linka odcinacza zaczepiona o zwalniacz. Odcinacz uzbrojony.}{}
\krok{Gondole zamknięte, owinięte taśmą, pasy napięte, opis markerem widoczny.}{}
\etap{Ostatnia kontrola}
\krok{Wejdź na \www{aprs.fi} i SondeHub: ostatnia pozycja ma mniej niż 2 minuty.}{\pole[2cm]{godz.:}}
\krok{SPOT i znaczniki Bluetooth włączone.}{}
\krok{Zadzwoń do służby ruchu lotniczego (jeśli wymagane w zgodzie).}{\pole[2cm]{godz.:}}
\krok{Kierownik startu: niebo wolne, widzowie w strefie, zespół na pozycjach.}{}
\krok{Kierownik misji: decyzja „start”.}{}
\etap{Start}
\krok{Osoba z gondolą ustawia się pod wiatr (po zawietrznej, twarzą do balonu).}{}
\krok{Odliczanie. PM zdejmują folię i puszczają balon.}{}
\krok{Balon nad gondolą i spokojny: komenda „start”. Osoba z gondolą ją puszcza.}{}
\krok{Zapisz czas startu.}{\pole[2cm]{UTC:}}
\krok{Fotograf: seria zdjęć / film startu.}{}
\end{procedura}
\komentarze[3]
\podpisy

% ===========================================================================
\formularz[operator śledzenia]{F12}{Lot i pościg}{T0}
Zapisuj pozycję co kilka minut. Wpisz zdarzenia: pęknięcie, odcięcie, utrata sygnału, zmiana trasy.

\renewcommand{\arraystretch}{1.75}
\footnotesize
\begin{tabularx}{\linewidth}{@{}|P{13mm}|P{16mm}|P{28mm}|P{28mm}|P{14mm}|L|@{}}\hline
\rowcolor{sdlight}\textbf{Godz.} & \textbf{Wys. [m]} & \textbf{Szer. geogr.} & \textbf{Dł. geogr.} & \textbf{Źródło} & \textbf{Zdarzenie, predykcja lądowania}\\\hline
 & & & & & \\\hline
 & & & & & \\\hline
 & & & & & \\\hline
 & & & & & \\\hline
 & & & & & \\\hline
 & & & & & \\\hline
 & & & & & \\\hline
 & & & & & \\\hline
 & & & & & \\\hline
 & & & & & \\\hline
 & & & & & \\\hline
 & & & & & \\\hline
 & & & & & \\\hline
 & & & & & \\\hline
 & & & & & \\\hline
 & & & & & \\\hline
 & & & & & \\\hline
\end{tabularx}
\normalsize
\podpisy

% ===========================================================================
\formularz[kierownik misji]{F13}{Odzyskanie}{T0}
\begin{procedura}
\etap{Dojazd}
\krok{Zapisz ostatnią pozycję z trackera.}{\pole[2cm]{wsp.:}}
\krok{Sprawdź pozycję SPOT i znaczników Bluetooth.}{}
\krok{Ustal właściciela terenu. Uzyskaj zgodę na wejście (pole, las prywatny).}{}
\etap{Na miejscu}
\krok{Fotograf: zdjęcia miejsca lądowania, zanim ktoś dotknie sprzętu.}{}
\krok{Sprawdź stan olinowania i odcinacza. Zapisz, czy odcinacz zadziałał.}{}
\krok{Wyłącz elektronikę platformy. Zabezpiecz karty pamięci.}{}
\krok{Przekaż ładunek opiekunowi ładunku.}{}
\krok{Zbierz resztki powłoki i spadochron. Nie zostawiaj śmieci.}{}
\etap{Zakończenie}
\krok{Zgłoś koniec lotu (jeśli wymagane).}{\pole[2cm]{godz.:}}
\krok{Sprawdź: wszyscy członkowie zespołu i cały sprzęt w samochodach.}{}
\end{procedura}
\komentarze[3]
\podpisy

% ===========================================================================
\formularz[kierownik misji]{F14}{Po misji}{do T+7 dni}
\begin{procedura}
\krok{Skopiuj dane z kart pamięci, trackera i serwisów (\www{aprs.fi}, SondeHub) do repozytorium.}{}
\krok{Zbierz zdjęcia od fotografa i zespołu w jednym miejscu.}{}
\krok{Zrób spotkanie podsumowujące (debrief) z całym zespołem.}{\pole[2cm]{data:}}
\krok{Napisz raport z misji (rozdział~\ref{ch:projekt}).}{}
\krok{Dopisz misję i wnioski do rozdziału~\ref{ch:historia}.}{}
\krok{Przenieś notatki terenowe i poprawki do podręcznika. Wydaj nową wersję.}{\pole[2cm]{wersja:}}
\krok{Dołącz dokumentację do wydania kodu lotnego.}{}
\krok{Sprawdź, napraw i odłóż sprzęt do szafki z opisem.}{}
\krok{Wyłącz subskrypcję SPOT. Rozlicz koszty.}{}
\end{procedura}
\noindent\textbf{Co działało dobrze:}\foreach \i in {1,...,2}{\linia}\par
\noindent\textbf{Co nie działało i dlaczego:}\foreach \i in {1,...,2}{\linia}\par
\noindent\textbf{Co zmienić w następnej misji:}\foreach \i in {1,...,2}{\linia}\par
\podpisy
